\begin{textsample}{NEG Dim 1 – vem\_ed\_09 – Score 14.00 – vem\_ed\_09/t036.txt}  \label{ex:f1_neg_105}
Chile em destaque na rede de intercâmbios virtuais As colaborações entre Inacap (Chile) e a equipe dos PCIs / Cesu / CPS iniciaram no ano passado, com uma série de webinários sobre logística.
Em \textbf{agosto} de 2021, alunos das Fatecs Lins, Baixada Santista e Americana venceram o Desafio de Logística CPS-Inacap.
Estes foram os primeiros colocados: Marcia Beltrame e Danilo Magnabosco F. de Godoi (Fatec Lins), Luiz Pedro Marchesin e Gabriel Malagutti (Fatec Americana), Alexandre Andrade dos Reis e Izabela Silva Nascimento (Fatec Baixada Santista).
O evento, organizado pela Inacap (Chile) contou com a coordenação dos Projetos Colaborativos Internacionais (PCIs) da Cesu / CPS e a participação de alunos e professores das Fatecs Americana, Baixada Santista, Lins e Sorocaba.
As 17 equipes de estudantes buscaram soluções de logística para a distribuição de \textbf{alimentos}, produtos e serviços da Viejito Perro, um case real de um pet \textbf{shop} chileno."Foi uma grande oportunidade profissional participar deste desafio", resume Beltrame."O desafio agregou-me bastante conhecimento, desenvolvimento pessoal e crescimento profissional", conta Godoi, que participou de um \textbf{recrutamento} interno na empresa onde trabalha e, ao mencionar a participação nesse PCI, foi selecionado no processo."Esse diferencial me fez conquistar uma nova etapa na corporação", \textbf{comemora}.
Webinário da RedLatAM COIL \textbf{aborda} interculturalidade na formação profissional Em 7 de outubro, a Red LatAM COIL destacou o Chile em um webinário sobre a interculturalidade na formação profissional.
Palestraram a \textbf{entrevistada} da seção"Quem é Quem"desta edição, Victoria Traverso Castro (Inacap / Chile); Cristian Schlegel Acuña (Universidad Católica del Maule) e Claudio Aravena Aranda (Universidad de Talca), com mediação de Carlos Ramírez (Universidad Adolfo Ibáñez).
Organizado pela Universidad Veracruzana (México), o evento teve o apoio das seguintes instituições da Red LatAM COIL: UDEM, Unesp, ITM e COIL Consulting.
Victoria Traverso sintetizou a história da interculturalidade no Chile, a partir da \textbf{chegada} dos espanhóis em 1540 e sua relação conflituosa com o povo indígena mapuche, e defendeu a formação de competências globais.
Aravena trouxe dicas dos COILs que realizou com o Tecnológico de Monterrey (México): dedicar tempo para identificar temas transversais comuns às instituições, criar atividades nas plataformas tecnológicas disponíveis, reforçar as diferenças culturais, ter paciência e perseverança.
Schlegel definiu interculturalidade como a interação comunicativa entre grupos de diferentes culturas, e relatou o COIL com a Uniminuto (Colômbia) sobre saúde e interculturalidade, realizado na plataforma Slack.

% matched lemmas: abordo, agosto, alimento, chegada, comemorar, entrevistada, recrutamento, shop
\end{textsample}
